\section{Computable spectral gap enclosures}\label{inf:sec-gap-enclosures}

This appendix supplies the optional selection rule of
Remark~\ref{inf:split-selection} and is not part of the proof
chain of Theorems~\ref{inf:theorem-A} and \ref{inf:theorem-B}.

Fix an input $p,R,\theta$ and the finite set of integer splits in a
compact interval of Section~\ref{inf:sec-widened}. Represent each
quartic Dirichlet operator $H_a$ on $\Hc$ by the harmonic chart
\[
 q=1+z_{\mathrm q}S_2,\qquad d=1+5z_{\mathrm q}S_2,\qquad
 z_{\mathrm q}=-\frac{\theta\mathsf H_2(1)}{2R^2pg_2}.
\]
All parameters of this chart are computable from an isolated
Neumann zero $R$ and the rational block parameters. The following
bound supplies a strict starting condition for the spectral enclosure.

\begin{lemma}[A computable quartic relative bound]
\label{inf:quartic-computable-relative}
Assume integer block sizes with $a+b=2p$, $\bb\ge\ab\ge16$,
$p\ge64$, and
$\mu_{\mathrm{alg}}=|z_{\mathrm q}|\le1/(100p)$.
Then the positive unitary quartic operator satisfies
\begin{equation}\label{inf:algorithm-relative-bound}
 \nrm{(H_a-H_0)H_0^{-1}}_{\Hc\to\Hc}
 \le\eta_{\mathrm{alg}},\qquad
 \eta_{\mathrm{alg}}=4000p\mu_{\mathrm{alg}}.
\end{equation}
For each fixed compact split interval, the condition
$\eta_{\mathrm{alg}}<1/2$ holds at every sufficiently large
input of either theorem, uniformly over its candidate splits.
\end{lemma}
\begin{proof}
Write $\mu=\mu_{\mathrm{alg}}$ within this proof.
The endpoint proof in Lemma~\ref{inf:angular-sup} uses
$\bb-1\ge\ab-1\ge0$ and therefore applies to these parameters.
Homogeneity and its derivative bounds, at degree four, give
\[
 \nrm{S_2}_\infty\le1,\qquad
 \nrm{\nabla S_2}_\infty\le4\sqrt2<6,
 \qquad \nrm{D^2S_2}_{\mathrm{op},\infty}\le32.
\]
For the last bound, at a unit point the radial-radial entry is
$12G_2$, the radial-tangential entries are $3\nabla_SG_2$, and
the tangential block is $\nabla_S^2G_2+4G_2g_S$.
The diagonal blocks have norm at most twenty and the off-diagonal
block at most twelve. Homogeneity gives the same bound inside the ball.

Thus $q,d>0$, $q^{-1},d^{-1}\le2$, and
\[
 \nrm{D\Phi^{-1}}_{\mathrm{op},\infty}\le3,\qquad
 \nrm{A_q}_{\mathrm{op},\infty}\le9,
 \qquad \nrm{A_q-I}_{\mathrm{op},\infty}\le200\mu.
\]
Indeed the inverse is $q^{-1}(I-y\otimes\nabla q/d)$.
In \eqref{inf:pullback-metric}, the last bound follows from
\[
 |q^{-2}-1|\le8\mu,\qquad
 q^{-2}\left(2|\nabla q|/d+|\nabla q|^2/d^2\right)
 \le96\mu+576\mu^2.
\]
The drift vector in the scalar coordinate formula is
$-D\Phi^{-1}(2A_q\nabla q+y(A_q:D^2q))$.
Its norm is at most
$3(108+288n)\mu\le2000p\mu$, using $n=2p$.
For $u=H_0^{-1}f$, the ball estimates give
$\nrm{D^2u}_2\le\nrm f_2$ and
$\nrm{\nabla u}_2\le(p-1)^{-1}\nrm f_2$.
The Frobenius norm of a symmetric $n$-matrix is at most $\sqrt n$
times its operator norm. Consequently, for $H_\Phi=-P_\Phi\Delta$,
\begin{equation}\label{inf:algorithm-scalar-bound}
 \nrm{(H_\Phi-H_0)H_0^{-1}}
 \le200\sqrt{2p}\,\mu+\frac{2000p\mu}{p-1}
 \le500p\mu.
\end{equation}

Put $m=q^{p-1/2}d^{1/2}$ and $l=\log m$.
Since $p\mu\le1/100$, the inequalities
$\log(1+t)\le t$ and $-\log(1-t)\le2t$ on $0\le t\le1/2$
give $\nrm m_\infty,\nrm{m^{-1}}_\infty\le2$.
Harmonicity of $q,d$ gives
\[
 \nrm{\nabla l}_\infty\le13p\mu,
 \qquad \nrm{\Delta l}_\infty\le200p\mu^2.
\]
For example the gradient bound follows from
$2(p-1/2)6\mu+5\cdot6\mu\le13p\mu$;
the Laplacian bound follows from
$4p(6\mu)^2+2(30\mu)^2\le200p\mu^2$.
The identity
$\Delta m^{-1}=m^{-1}(|\nabla l|^2-\Delta l)$ therefore yields
\[
 \nrm{[H_0,m^{-1}]H_0^{-1}}
 \le2\left\{\frac{26p\mu}{p-1}
      +\frac{169p^2\mu^2+200p\mu^2}{(p-1)^2}\right\}
 \le60\mu.
\]
Here $p\ge64$ and $\mu\le1/(100p)$ bound the braces by $30\mu$.
It follows that $\nrm{H_0m^{-1}H_0^{-1}}\le3$.
Finally use $H_a=mH_\Phi m^{-1}$ and the exact factorization
\begin{align*}
 (H_a-H_0)H_0^{-1}
 ={}&m(H_\Phi-H_0)H_0^{-1}(H_0m^{-1}H_0^{-1})\\
    &+m[H_0,m^{-1}]H_0^{-1}.
\end{align*}
Its norm is at most $2(1500p\mu+60\mu)\le4000p\mu$.
All estimates hold on the Dirichlet graph core and extend by its
density. In the families under consideration,
$\mu_{\mathrm{alg}}\le C_{\mathrm{sp}}\tau p^{-2}$, hence
$\eta_{\mathrm{alg}}\le C_{\mathrm{sp}}\tau/p\to0$ uniformly for
fixed detuning bounds and for $\tau\le\log p$.
\end{proof}

\begin{lemma}[Convergent gap enclosures]\label{inf:gap-enclosures}
For a computable input satisfying the hypotheses of
Lemma~\ref{inf:quartic-computable-relative} and
$\eta_{\mathrm{alg}}<1/2$, the capped gap
\[
 g_a=\min\{R^2,\dist(R^2,\spec H_a)\}
\]
has rational lower and upper enclosures of arbitrarily small
prescribed width.
\end{lemma}
\begin{proof}
Set $\mathcal R_{\mathrm{alg},a}=(H_a-H_0)H_0^{-1}$ and
$\mathcal K_a=H_a^{-1}$. The strict relative bound gives
\begin{equation}\label{inf:compact-inverse-series}
 \mathcal K_a=H_0^{-1}\sum_{v=0}^{\infty}(-\mathcal R_{\mathrm{alg},a})^v,
 \qquad
 \left\|H_0^{-1}\sum_{v>L}(-\mathcal R_{\mathrm{alg},a})^v\right\|
 \le\frac{\eta_{\mathrm{alg}}^{L+1}}
 {(p-1)^2(1-\eta_{\mathrm{alg}})}.
\end{equation}
The operator $\mathcal K_a$ is compact, positive and self-adjoint.
Use the orthonormal Bessel--Jacobi basis of $H_0$ and the rectangular
projection $\Pi^{\mathrm{alg}}_{J,K_{\mathrm{rad}}}$ onto
$0\le j\le J$, $1\le k\le K_{\mathrm{rad}}$.
The radial min--max principle gives
\[
 z_{j,k}^2\ge\nu_j^2-1/4+\pi^2 k^2.
\]
Thus one may take the omitted-eigenvalue lower bound
\[
 \Lambda_{\mathrm{miss}}=\min\{\nu_{J+1}^2-1/4+\pi^2,
                   \ \nu_0^2-1/4+\pi^2(K_{\mathrm{rad}}+1)^2\}.
\]
It tends to infinity as $J,K_{\mathrm{rad}}\to\infty$.
For this projection, abbreviated $\Pi^{\mathrm{alg}}$, commutation
with $H_0^{-1}$ and self-adjointness give
\begin{equation}\label{inf:compact-inverse-tail}
 \nrm{\mathcal K_a-\Pi^{\mathrm{alg}}\mathcal K_a\Pi^{\mathrm{alg}}}
 \le\frac{2}{(1-\eta_{\mathrm{alg}})\Lambda_{\mathrm{miss}}}.
\end{equation}
Indeed $(I-\Pi^{\mathrm{alg}})\mathcal K_a$ is bounded by the
right side without its factor two, using
$\mathcal K_a=H_0^{-1}(I+\mathcal R_{\mathrm{alg},a})^{-1}$.
Its adjoint has the same bound, and their sum gives
\eqref{inf:compact-inverse-tail}.

We describe the finite matrix computations in these two truncations.
The basis values and derivatives are computable from the convergent
Bessel series, finite Jacobi polynomials and their normalizations.
The positive roots are simple and separated by more than $\pi$.
The sine comparison in Lemma~\ref{inf:radial-poles}, applied on
disjoint length-four intervals beyond $2\nu$, gives the coarse
bound $j_{\nu,k}\le2\nu+8k$. Since the first root exceeds
$\nu>1$, a rational mesh starting at one, with gaps below $\pi$,
finds every root up to a prescribed index by sign changes.
Mesh points can be chosen with certified nonzero values: of two
candidate points less than $\pi$ apart, at least one is not a zero,
and successive series enclosures detect that value. Choosing them
in small disjoint neighbourhoods of an initial mesh preserves
its order and gap bound. The same procedure refines each isolated
root interval. Simplicity then gives a nonzero derivative bound
after finitely many refinements, supplying its normalization.

For a finite basis sum $v$, the function
$\mathcal R_{\mathrm{alg},a}v$ is an explicit smooth invariant
function. Its coefficients and its squared norm are integrals in
the radial and angular variables with positive weights.
Every chart denominator is bounded away from zero by
$q\ge1-\mu_{\mathrm{alg}}$, $d\ge1-5\mu_{\mathrm{alg}}$.
The regular Bessel series times solid harmonics gives computable
uniform bounds for the function and its required derivatives on
the closed ball: each derivative series has factorial tails and
only finitely many input modes occur.
After the coordinate substitution, use an endpoint cutoff
$0<\epsilon_{\mathrm{cut}}<1/2$. The normalized radial and angular
densities are $2p\zeta^{2p-1}$ and
$x^{\ab-1}(1-x)^{\bb-1}/B(\ab,\bb)$. Their omitted endpoint
measures are bounded, respectively, by
\[
 \frac{2p\epsilon_{\mathrm{cut}}^{2p}}{2p},\qquad
 \frac{\epsilon_{\mathrm{cut}}^{\ab}}{\ab B(\ab,\bb)},\qquad
 \frac{\epsilon_{\mathrm{cut}}^{\bb}}{\bb B(\ab,\bb)}
\]
at the origin and the two angular endpoints. Multiplying by the
computed supremum bounds controls those tails.
On the remaining compact rectangle, derivative bounds give a
computable modulus of continuity and rational interval Riemann
sums converge from above and below. Here
$B(\ab,\bb)=\Gamma(\ab)\Gamma(\bb)/\Gamma(p)$.
For these integer or half-integer parameters the Gamma formulas
give convergent rational enclosures; a positive lower enclosure
of $B(\ab,\bb)$ gives the required upper bound for its reciprocal.
Equivalently one can obtain this lower enclosure from its positive
integral. Thus the tail constants $2p$ and $1/B(\ab,\bb)$ have
computable rational upper bounds. This constructs convergent
rational enclosures of every required integral.

To approximate the image of a finite vector, use Parseval in this
complete basis. For a larger finite projection $\Pi'$, its squared
residual is
\[
 \nrm{\mathcal R_{\mathrm{alg},a}v}^2
 -\sum_{i\in\Pi'}
       |\langle\mathcal R_{\mathrm{alg},a}v,e_i\rangle|^2.
\]
Its upper enclosure eventually falls below any positive target by
completeness, increasing the projection and integration precision.
This supplies a terminating precision search for each application
of $\mathcal R_{\mathrm{alg},a}$. An input error is multiplied by at
most $\eta_{\mathrm{alg}}$ and added to the next approximation error.
There are only finitely many applications in the truncated series
\eqref{inf:compact-inverse-series}. Taking errors sufficiently small
for its finitely many input columns gives a rational approximation
to the finite compression in operator norm. A column error at most
the requested tolerance divided by the number of columns suffices,
by Cauchy--Schwarz. Symmetrization preserves
that error bound.

The spectrum of a finite rational symmetric matrix can be enclosed
by rational arithmetic. One construction enumerates rational
orthogonal matrices formed from coordinate reflections and plane
rotations with entries
$(1-t^2)/(1+t^2)$ and $2t/(1+t^2)$, $t\in\mathbb Q$.
This family is dense in the orthogonal group. Continue until the
off-diagonal Frobenius norm of the conjugated matrix is below a
prescribed tolerance. The spectral theorem guarantees termination,
and the diagonal entries then enclose its spectrum with that error.
For self-adjoint operators an operator-norm error bounds their
spectral Hausdorff distance: outside that neighbourhood the
resolvent identity and a Neumann series give invertibility, in both
directions. Combining this fact with
\eqref{inf:compact-inverse-series}--\eqref{inf:compact-inverse-tail}
gives arbitrary norm-accurate finite spectral enclosures of
$\mathcal K_a$, with zero adjoined for the infinite complement.

Finally put
\[
 f_R(t)=\min\{R^2,|t^{-1}-R^2|\}\quad(t>0),\qquad
 f_R(t)=R^2\quad(t\le0).
\]
It is continuous and Lipschitz with constant $4R^4$: its nonconstant
part has $t\ge1/(2R^2)$ and derivative of absolute value at most
$1/t^2$. Also it is Lipschitz with constant one as a function of
$R^2$, so the input radius can be enclosed at the same time.
Spectral inversion gives
$g_a=\min f_R(\spec\mathcal K_a)$.
The preceding finite spectral enclosures therefore give rational
bounds on $g_a$ with error tending to zero. Clipping the lower bound
at zero preserves the enclosure. This proves the assertion.
\end{proof}
